% SSSUP-FONTE-STORICA: fonti-locali/originali/2014-2020.zip
% Anno 2016, esercizio 1. Testo storico; figura ricostruita da PDF.
% Verificare editorialmente la trasposizione prima di considerarla definitiva.
\renewcommand{\Anno}{2016}
\renewcommand{\NumeroProblema}{1}

\renewcommand{\Testo}{
Una rastrelliera di prese è costituita da una base di plastica rettangolare molto lunga e di larghezza $l$. Al centro della base sono fissate $n$ prese a distanza $h$ l'una dall'altra. Una rastrelliera di prese $P$ ed una di carichi $C$ vengono fissate, longitudinalmente e mediante una guida che corrisponde al loro asse, a due binari che distano $l$ tra loro e lasciate libere di scorrere verticalmente.

\begin{center}
% SSSUP-FIGURA-RICOSTRUITA: schema rastrelliere; fonte 2016figura.pdf.
\psset{unit=0.88cm}
\begin{pspicture}(-0.6,-0.7)(6.8,4.7)
  \psframe[fillstyle=solid,fillcolor=lightgray,linewidth=.75pt](1.0,0.45)(2.1,3.65)
  \psframe[fillstyle=solid,fillcolor=white,linewidth=.75pt](3.55,1.05)(4.65,4.25)
  \psline[linestyle=dashed,linewidth=.45pt](1.55,0.0)(1.55,4.7)
  \psline[linestyle=dashed,linewidth=.45pt](4.10,0.0)(4.10,4.7)
  \rput[r](.82,2.3){$P$}
  \rput[l](4.83,3.6){$C$}
  \psline[linestyle=dotted](2.10,4.4)(3.55,4.4)
  \psline{<->}(2.10,4.35)(3.55,4.35)
  \rput[b](2.82,4.45){$\ell$}
  \psline{<->}(4.84,2.65)(4.84,3.02)
  \rput[l](5.03,2.86){$h$}
  \pscircle*[linecolor=black](1.55,0.78){.080}
  \pscircle*[linecolor=black](4.10,1.33){.080}
  \pscircle*[linecolor=black](1.55,1.17){.080}
  \pscircle*[linecolor=black](4.10,1.72){.080}
  \pscircle*[linecolor=black](1.55,1.56){.080}
  \pscircle*[linecolor=black](4.10,2.11){.080}
  \pscircle*[linecolor=black](1.55,1.95){.080}
  \pscircle*[linecolor=black](4.10,2.50){.080}
  \pscircle*[linecolor=black](1.55,2.34){.080}
  \pscircle*[linecolor=black](4.10,2.89){.080}
  \pscircle*[linecolor=black](1.55,2.73){.080}
  \pscircle*[linecolor=black](4.10,3.28){.080}
  \pscircle*[linecolor=black](1.55,3.12){.080}
  \pscircle*[linecolor=black](4.10,3.67){.080}
\end{pspicture}
\end{center}
Si desidera connettere tutte le prese della rastrelliera $C$ dei carichi ad alcune (non necessariamente tutte) di quelle della rastrelliera $P$ tenendo conto dei seguenti fattori:
\begin{itemize}
\item il filo per le connessioni costa $\beta$ per unità di lunghezza;
\item ciascuna presa di $P$ sopporta al massimo il carico di 2 prese di $C$;
\item uno sdoppiatore è puntiforme, non costa niente e può essere inserito in qualsiasi punto del cavo di connessione;
\item connettersi ad una presa di $P$ costa $\alpha=\beta/3$.
\end{itemize}

\begin{enumerate}
\item[(a)] Per $n$ pari, si decida, al variare di $l,h,\beta$ quali sono le configurazioni di connessione più economiche.
\item[(b)] Si analizzi, in modo qualitativo se non fosse possibile calcolare con esattezza le soglie, il caso $n$ dispari.
\end{enumerate}
}

\renewcommand{\Risposta}{}
\renewcommand{\Aiutino}{}
\renewcommand{\SoluzioneUfficiale}{}
\renewcommand{\SoluzioneRedazionale}{}
\renewcommand{\SoluzioniAlternative}{}
\renewcommand{\Osservazioni}{}
\renewcommand{\Commenti}{}
